\scanpage{72}
% archive foliation: 71
% Typescript is upside down in the supplied whole-page orientation.
\deleted{
Ceci donne une démonstration plus sympatique et élémentaire de 11.9.6. Quant à 11.9.3. (ii), il varésulter d'un

\uline{Lemme} (ou corollaire à la proposition précédente). Sous les hypothèses de la proposition, ne supposons plus néc que $A'$ soit projectif sur $A$, mais seulement ceci: pour tout quotient intègre $B=A/\underline{p}$ de $A$, posant $B'=B\otimes_A A'$, il existe un $h\ B-(0)$ tel que $B_h$ soit \deleted{projectif} \inserted{libre} sur $A_h$.
Alors la conclusion de la proposition reste encore valable.

Soit en effet $\phi$ une section de $F'$ dont les images par les $u'$ sont nulles, prouvons $\phi=0$, i.e. $\mathrm{supp}\,\phi=\varnothing$. Dans le cas contraire, il existerait un point maximal $x'$ dans $\mathrm{supp}\,\phi$, correspondant à l'idéal premier $\underline{p}'$ de $A'$. Soit \deleted{$\underline{p}=\unsure{\underline{p}'\cap A'|\underline{p}'\cap A}$} \inserted{$x$} son image dans $X$, défini par l'idéal $\underline{p}$ trace de $\underline{p}'$ sur $A$, alors en vertu de l'hypothèse appliquée à $A/\underline{p}$, on voit que, quitte à remplacer $X$ par un voisinage ouvert affine de $x$, on peut supposer que $\unsure{A'|B'}/\underline{p}B'$ est \deleted{projectif} \inserted{libre} sur $A/\underline{p}$. Notons que nous pouvions choisir $x'$ de façon que $x$ soit également maximal dans $g(\mathrm{supp}\,\phi)$ (commencer par choisir $x$ maximal dans $g(\mathrm{supp}\,\phi)$, puis $x'$ maixmal dans $g^{-1}(x)$), ce qui signifie que localisant en $\underline{p}$, le support de $\phi_{\underline{p}}$ est contenu dans $\mathrm{Spec}(A'_{\underline{p}}/\underline{p}A'_{\underline{p}})$, donc $\phi_{\underline{p}}$ est annulé par une puissance de $\underline{p}\,A'_{\underline{p}}$, ou encore que l'on peut trouver $h\ A-\underline{p}$ tel que
}
